\chapter{Introduction}
\label{ch:introduction}

Rate limiting protects a shared resource by rejecting or delaying requests once a client
exceeds an agreed budget. The textbook implementation keeps a counter behind a lock: every
request checks the counter, and the counter lives on one machine so that every client sees
the same number. That design works until the service the rate limiter protects itself
becomes distributed, at which point the counter becomes a second service that every request
must reach before it can reach the first one. The coordinator meant to bound latency has
become a new source of it.

This thesis asks whether a rate limiter can give the same guarantee---no client exceeds its
budget by more than a bounded, known amount---without a machine that every request must
consult. We answer yes, for the common case where clients are willing to accept a
probabilistic bound instead of an exact one, using a construction we call \sysname{}.
\sysname{} lets each node estimate the global request rate from a small amount of gossiped
state and make an admit-or-reject decision locally, in the time it takes to read a counter
from memory.

\section{Contributions}
\label{sec:intro-contributions}

The thesis makes three contributions, in the order a reader would want to evaluate them.
First, we give a \term{gossip-based} estimator for the global request rate that converges
within a bounded number of rounds regardless of cluster size, described in
Section~\ref{sec:token-bucket}. Second, we show that combining this estimator with a local
\term{token bucket} preserves the standard token-bucket guarantee up to an error term that
shrinks as the gossip fan-out grows. Third, we evaluate the construction against a
centralized baseline on a synthetic workload and report the latency and accuracy trade-off
in Chapter~\ref{ch:results}; the headline result is a reduction in tail latency with no
client observing more than a small, bounded excess over its budget.

\section{Why a coordinator is the wrong default}

The appeal of a central counter is that it is trivially correct: there is exactly one
number, and every decision reads and writes it under a lock. The cost is that the counter
sits on the path of every request the system serves, so its own latency and availability
become the rate limiter's latency and availability. A rate limiter that is slower or less
reliable than the service it protects has made things worse, not better, and this failure
mode is easy to miss in a single-region deployment where the coordinator's round trip is a
few milliseconds. It stops being easy to miss the moment the deployment crosses a region
boundary, where the same round trip is measured in tens of milliseconds and the coordinator
becomes the dominant term in every request's latency budget.
